\begin{ouverture}[calcul]{6}{Équations et résolution de problèmes}{Partie 2 · Algèbre et modèles}
{Inconnue · Équilibre · Mise en équation · Vérification}
\objectif{Tester si une valeur est solution.}
\objectif{Résoudre une équation du premier degré.}
\objectif{Traduire un problème puis interpréter la solution.}
\prerequis{Calcul littéral et nombres relatifs.}
\end{ouverture}
\pagestyle{otMain}\essentiel\fontsize{11}{15}\selectfont
\phantomsection\label{cours-6-1}
\essSection{1}{Conserver l’égalité}
\begin{definition}Une équation est une égalité contenant une inconnue. Une solution est une valeur qui rend cette égalité vraie. Résoudre une équation consiste à trouver toutes ses solutions dans l’ensemble de nombres considéré.\end{definition}
\begin{propriete}Ajouter ou soustraire un même nombre aux deux membres conserve les solutions. Multiplier ou diviser les deux membres par un même nombre non nul conserve aussi les solutions. Ces transformations permettent d’isoler l’inconnue.\end{propriete}
\begin{exemple}$3x+5=17$. On soustrait 5 aux deux membres : $3x=12$. On divise par 3 : $x=4$. Contrôle : $3\times4+5=17$. La solution est bien 4.\end{exemple}
\label{fig-cours-6-1-cours-equilibre}\figureGeo{cours-equilibre}

\begin{attention}« Passer de l’autre côté » doit correspondre à une opération faite aux deux membres. On ne change pas arbitrairement un signe ou une multiplication en addition.\end{attention}
\clearpage
\phantomsection\label{cours-6-2}
\essSection{2}{Du problème à la réponse}
\begin{savoirfaire}[Une méthode en quatre temps]1. Choisir l’inconnue et son unité. 2. Traduire la relation par une équation. 3. Résoudre en écrivant des égalités équivalentes. 4. Contrôler dans le contexte et formuler une phrase de réponse.\end{savoirfaire}
\begin{exemple}$5x-4=2x+8$ donne $3x-4=8$, puis $3x=12$ et $x=4$. Pour $2(x+3)=x+11$, développer donne $2x+6=x+11$, donc $x=5$.\end{exemple}
\begin{savoirfaire}[Plusieurs termes]Développer les parenthèses, réduire chaque membre, réunir les termes en $x$ d’un côté et les constantes de l’autre. Avec des fractions simples, multiplier les deux membres par un dénominateur commun non nul peut alléger les calculs.\end{savoirfaire}
\begin{attention}Une solution calculée doit respecter le problème : longueur positive, effectif entier, temps possible. Une égalité finale impossible comme $0=3$ signifie qu’il n’y a pas de solution. Une égalité toujours vraie comme $0=0$ signifie que toutes les valeurs admises conviennent.\end{attention}

\clearpage\rubrique{Découverte}
\begin{blocLivre}{ex-6-1}\begin{exercice}[1]{Tester une valeur}
Le nombre 4 est-il solution de $3x-2=10$ ? Le nombre 5 l’est-il ?
\end{exercice}\end{blocLivre}
\begin{blocLivre}{ex-6-2}\begin{exercice}[1]{Une addition}
Résoudre $x+8=-3$.
\end{exercice}\end{blocLivre}
\begin{blocLivre}{ex-6-3}\begin{exercice}[1]{Un coefficient négatif}
Résoudre $-4x=18$.
\end{exercice}\end{blocLivre}
\begin{blocLivre}{ex-6-4}\begin{exercice}[1]{Deux étapes}
Résoudre $5x-7=23$.
\end{exercice}\end{blocLivre}
\begin{blocLivre}{ex-6-5}\begin{exercice}[1]{Traduire une phrase}
Le triple d’un nombre augmenté de 4 vaut 25. Écrire l’équation puis trouver le nombre.
\end{exercice}\end{blocLivre}

\rubrique{Entraînement}
\begin{blocLivre}{ex-6-6}\begin{exercice}[2]{L’inconnue des deux côtés}
Résoudre $7x-5=3x+11$.
\end{exercice}\end{blocLivre}
\begin{blocLivre}{ex-6-7}\begin{exercice}[2]{Avec une parenthèse}
Résoudre $3(x-2)=x+8$.
\end{exercice}\end{blocLivre}
\begin{blocLivre}{ex-6-8}\begin{exercice}[2]{Avec des fractions}
Résoudre $\frac{x}{3}+\frac12=\frac52$.
\end{exercice}\end{blocLivre}
\begin{blocLivre}{ex-6-9}\begin{exercice}[2]{Corriger une méthode}
Un élève résout $2x+6=14$ en écrivant $x+6=7$, puis $x=1$. Corriger.
\end{exercice}\end{blocLivre}
\begin{blocLivre}{ex-6-10}\begin{exercice}[2]{Aucune ou toutes}
Résoudre $4x+2=4x-1$, puis $3(x+2)=3x+6$.
\end{exercice}\end{blocLivre}

\rubrique{Problèmes}
\begin{blocLivre}{ex-6-11}\begin{exercice}[3]{Un rectangle}
Un rectangle a une largeur $x$ cm et une longueur $x+5$ cm. Son périmètre est 38 cm. Trouver ses dimensions.
\label{fig-ex-6-11}\figureGeo{eq-rectangle}
\end{exercice}\end{blocLivre}
\begin{blocLivre}{ex-6-12}\begin{exercice}[3]{Deux tarifs}
Le tarif A coûte 12 euros puis 3 euros par entrée. Le tarif B coûte 5 euros par entrée. Pour combien d’entrées les dépenses sont-elles égales ?
\end{exercice}\end{blocLivre}
\begin{blocLivre}{ex-6-13}\begin{exercice}[3]{Des entiers consécutifs}
La somme de trois entiers consécutifs vaut 81. Les trouver.
\end{exercice}\end{blocLivre}
\begin{blocLivre}{ex-6-14}\begin{exercice}[3]{Un âge futur}
Maya a 12 ans et son père 40 ans. Dans combien d’années l’âge du père sera-t-il le double de celui de Maya ?
\end{exercice}\end{blocLivre}
\begin{blocLivre}{ex-6-15}\begin{exercice}[3]{Une longueur admissible}
Un triangle équilatéral de côté $x+2$ cm et un carré de côté $x$ cm ont le même périmètre. Trouver $x$, puis vérifier que les figures existent.
\label{fig-ex-6-15}\figureGeo{eq-deux-figures}
\end{exercice}\end{blocLivre}

\clearpage
